\documentclass[11pt,a4paper]{article}

\usepackage[utf8]{inputenc}
\usepackage[main=italian,provide=*]{babel}
\usepackage{amsmath}
\usepackage{amssymb}
\usepackage{mathtools}
\usepackage{geometry}
\usepackage{booktabs}
\usepackage{hyperref}

\geometry{margin=2.7cm}

\title{Quantum simulation del dimero: derivazione teorica completa\\
\large Decomposizione di Suzuki--Trotter per l'Hamiltoniana con termine DM}
\author{}
\date{}

\newcommand{\ket}[1]{\lvert #1 \rangle}
\newcommand{\bra}[1]{\langle #1 \rvert}

\begin{document}

\maketitle
\tableofcontents

\bigskip
\noindent Questo documento raccoglie in forma sistematica le derivazioni relative al task di \emph{quantum simulation} (evoluzione temporale reale $e^{-iHt}$) richiesto dal relatore sull'Hamiltoniana VQE del dimero, in vista dell'implementazione tramite decomposizione di Suzuki--Trotter. Ogni passaggio è riportato per esteso, senza salti logici.

\section{Hamiltoniana e struttura generale}

L'Hamiltoniana del dimero, in convenzione $s_i^\alpha = \sigma_i^\alpha/2$, è

\begin{equation}
H = b(s_{z1}+s_{z2}) + J\,\vec s_1\cdot\vec s_2 + D(s_{x1}s_{z2}-s_{z1}s_{x2}) = H_1 + H_2,
\end{equation}

dove si è scelto di raggruppare i termini in due blocchi:

\begin{align}
H_1 &= b(s_{z1}+s_{z2}) + J\,\vec s_1\cdot\vec s_2 \qquad \text{(campo + scambio isotropo, \emph{fisso}, indipendente da } D\text{)},\\
H_2 &= D(s_{x1}s_{z2}-s_{z1}s_{x2}) \qquad \text{(solo termine DM, si annulla per } D=0\text{)}.
\end{align}

In generale, per $D\neq0$, i due blocchi non commutano: $[H_1,H_2]\neq 0$. Per $D=0$ si ha semplicemente $H=H_1$ e l'evoluzione temporale è \textbf{esatta}, senza alcuna approssimazione. Per $D\neq0$ è necessaria la decomposizione di Suzuki--Trotter, il cui errore verrà derivato nella Sezione~\ref{sec:trotter-error}.

Lo scopo delle Sezioni~\ref{sec:UJ}--\ref{sec:U2} è mostrare che, sebbene $H_1$ e $H_2$ non commutino tra loro, ciascuno dei due blocchi si decompone \emph{esattamente} (senza errore di approssimazione) in un prodotto di rotazioni a uno e due qubit implementabili nativamente in Qiskit.

\section{Decomposizione esatta di $U_J(t) = e^{-iJt\,\vec s_1\cdot\vec s_2}$}
\label{sec:UJ}

\subsection{Riscrittura in operatori di Pauli}

Il prodotto scalare degli spin si scrive esplicitamente come

\begin{equation}
\vec s_1\cdot\vec s_2 = s_{x1}s_{x2} + s_{y1}s_{y2} + s_{z1}s_{z2}.
\end{equation}

Sostituendo $s_i^\alpha = \sigma_i^\alpha/2$ in ciascun termine:

\begin{equation}
\vec s_1\cdot\vec s_2 = \frac{X_1X_2}{4} + \frac{Y_1Y_2}{4} + \frac{Z_1Z_2}{4} = \frac{1}{4}\left(X_1X_2 + Y_1Y_2 + Z_1Z_2\right),
\end{equation}

dove si è abbreviato $X_1X_2 \equiv X\otimes X$ (analogamente per $Y,Z$). Il termine di scambio isotropo diventa quindi

\begin{equation}
J\,\vec s_1\cdot\vec s_2 = \frac{J}{4}\left(XX+YY+ZZ\right),
\end{equation}

e l'operatore di evoluzione corrispondente è

\begin{equation}
U_J(t) = \exp\!\left[-i\,\frac{Jt}{4}\left(XX+YY+ZZ\right)\right].
\end{equation}

Il fattore $1/4$ emerge automaticamente dal passaggio da spin $\vec s_i=\vec\sigma_i/2$ a Pauli: è lo stesso fattore-4 tracciato altrove nel progetto tra la convenzione $J_{\mathrm{code}}$ (che moltiplica direttamente gli operatori di Pauli) e $J_{\mathrm{Crippa}}$ (che moltiplica gli operatori di spin $\vec s_i$).

\subsection{Commutatività di $XX$, $YY$, $ZZ$}
\label{subsec:commXYZ}

Si dimostra che i tre operatori a due qubit $XX$, $YY$, $ZZ$ commutano a due a due. La dimostrazione usa la proprietà del prodotto misto tra operatori su spazi tensoriali,

\begin{equation}
(A\otimes B)(C\otimes D) = (AC)\otimes(BD),
\end{equation}

e le relazioni di anticommutazione/prodotto delle matrici di Pauli su un singolo qubit:

\begin{equation}
XY = iZ,\quad YX=-iZ,\quad YZ=iX,\quad ZY=-iX,\quad ZX=iY,\quad XZ=-iY.
\end{equation}

\paragraph{Coppia $(XX,YY)$.}

\begin{align}
(X\otimes X)(Y\otimes Y) &= (XY)\otimes(XY) = (iZ)\otimes(iZ) = i^2\,(Z\otimes Z) = -(Z\otimes Z),\\
(Y\otimes Y)(X\otimes X) &= (YX)\otimes(YX) = (-iZ)\otimes(-iZ) = (-i)^2\,(Z\otimes Z) = -(Z\otimes Z).
\end{align}

I due prodotti coincidono: $(XX)(YY) = (YY)(XX) = -ZZ$, dunque $[XX,YY]=0$.

\paragraph{Coppia $(YY,ZZ)$.}

\begin{align}
(Y\otimes Y)(Z\otimes Z) &= (YZ)\otimes(YZ) = (iX)\otimes(iX) = -(X\otimes X),\\
(Z\otimes Z)(Y\otimes Y) &= (ZY)\otimes(ZY) = (-iX)\otimes(-iX) = -(X\otimes X).
\end{align}

I due prodotti coincidono: $(YY)(ZZ)=(ZZ)(YY)=-XX$, dunque $[YY,ZZ]=0$.

\paragraph{Coppia $(ZZ,XX)$.}

\begin{align}
(Z\otimes Z)(X\otimes X) &= (ZX)\otimes(ZX) = (iY)\otimes(iY) = -(Y\otimes Y),\\
(X\otimes X)(Z\otimes Z) &= (XZ)\otimes(XZ) = (-iY)\otimes(-iY) = -(Y\otimes Y).
\end{align}

I due prodotti coincidono: $(ZZ)(XX)=(XX)(ZZ)=-YY$, dunque $[ZZ,XX]=0$.

\medskip
\noindent Essendo $[XX,YY]=[YY,ZZ]=[ZZ,XX]=0$, per l'argomento generale dimostrato nella Sezione~\ref{sec:BCH} l'esponenziale della somma fattorizza \textbf{esattamente} (in qualunque ordine dei tre fattori, senza alcun errore di troncamento):

\begin{equation}
U_J(t) = \exp\!\left(-i\tfrac{Jt}{4}XX\right)\exp\!\left(-i\tfrac{Jt}{4}YY\right)\exp\!\left(-i\tfrac{Jt}{4}ZZ\right).
\end{equation}

\subsection{Mappatura sui gate nativi Qiskit}

Qiskit definisce, nella convenzione standard (\texttt{qiskit.circuit.library}, verificata numericamente su Qiskit 2.x tramite \texttt{RXXGate(0.7).to\_matrix()}):

\begin{equation}
R_{XX}(\theta) = e^{-i\frac{\theta}{2}XX}, \qquad R_{YY}(\theta) = e^{-i\frac{\theta}{2}YY}, \qquad R_{ZZ}(\theta) = e^{-i\frac{\theta}{2}ZZ}.
\end{equation}

La verifica numerica diretta per $\theta=0.7$ dà, sulla diagonale, $\cos(\theta/2)=\cos(0.35)=0.9394$ e, fuori diagonale (nel pattern proprio di $X\otimes X$), $-i\sin(\theta/2)=-i\sin(0.35)=-0.3429\,i$, in pieno accordo con la formula attesa.

Confrontando l'esponente della definizione Qiskit, $\theta/2$, con l'esponente ottenuto nella Sezione~\ref{sec:UJ} per il singolo fattore, $Jt/4$, si ricava

\begin{equation}
\frac{\theta}{2} = \frac{Jt}{4} \quad\Longrightarrow\quad \theta = \frac{Jt}{2}.
\end{equation}

Si ottiene quindi il risultato finale per il blocco di scambio:

\begin{equation}
\boxed{U_J(t) = R_{XX}\!\left(\frac{Jt}{2}\right)\,R_{YY}\!\left(\frac{Jt}{2}\right)\,R_{ZZ}\!\left(\frac{Jt}{2}\right)}
\end{equation}

\section{Evoluzione completa di $H_1$ (campo + scambio isotropo)}
\label{sec:H1}

\subsection{Il campo commuta con lo scambio isotropo}

Si introduce lo spin totale $\vec S = \vec s_1+\vec s_2$ e il relativo Casimir $S^2=\vec S\cdot\vec S$. Espandendo esplicitamente:

\begin{equation}
S^2 = (\vec s_1+\vec s_2)\cdot(\vec s_1+\vec s_2) = s_1^2 + s_2^2 + 2\,\vec s_1\cdot\vec s_2,
\end{equation}

da cui, isolando il prodotto scalare:

\begin{equation}
\vec s_1\cdot\vec s_2 = \frac{1}{2}\left(S^2 - s_1^2 - s_2^2\right).
\label{eq:s1s2-casimir}
\end{equation}

Per un sistema di due spin-$1/2$ fissati, $s_1^2$ e $s_2^2$ sono costanti (autovalori fissi dell'operatore di Casimir di singolo spin), quindi il termine di scambio isotropo è, a meno di una costante additiva, proporzionale al Casimir $S^2$ dello spin totale.

Il punto cruciale è che $S^2$, essendo il Casimir dell'algebra del momento angolare, è per costruzione uno \emph{scalare invariante sotto rotazioni}, mentre $S_z$ (componente $z$ dello spin totale, che compare nel termine di campo) è il \emph{generatore delle rotazioni attorno all'asse $z$}. Un operatore invariante per rotazioni commuta sempre con qualunque generatore di rotazione: è esattamente lo stesso motivo per cui, per un singolo spin, vale $[\vec s^{\,2},s_z]=0$. Formalmente:

\begin{equation}
[S^2, S_z] = 0 \qquad \text{(proprietà generale dell'algebra del momento angolare)}.
\end{equation}

Applicando questo risultato all'Eq.~\eqref{eq:s1s2-casimir}, con $S_z=s_{z1}+s_{z2}$:

\begin{align}
\left[J\,\vec s_1\cdot\vec s_2,\; b(s_{z1}+s_{z2})\right] &= \left[\frac{J}{2}\left(S^2-s_1^2-s_2^2\right),\, bS_z\right]\\
&= \frac{Jb}{2}\underbrace{[S^2,S_z]}_{=0} - \frac{Jb}{2}\underbrace{[s_1^2+s_2^2,\,S_z]}_{=0\text{ (costanti)}}\\
&= 0.
\end{align}

Si noti che questo argomento è concettualmente distinto (e più diretto) rispetto a un eventuale argomento basato sulla simmetria di scambio dei siti $1\leftrightarrow 2$: quest'ultima garantisce che $\vec s_1\cdot\vec s_2$ sia simmetrico sotto lo scambio delle particelle, ma non fornisce di per sé alcuna informazione sulla relazione di commutazione con $S_z$. La commutazione con il campo discende invece esclusivamente dal fatto che $S^2$ è il Casimir dell'algebra $\mathfrak{su}(2)$ dello spin totale.

\subsection{Fattorizzazione esatta di $e^{-iH_1t}$}

Poiché il campo e lo scambio isotropo commutano, per l'argomento generale della Sezione~\ref{sec:BCH} l'esponenziale della loro somma fattorizza esattamente:

\begin{equation}
e^{-iH_1t} = e^{-i\left[b(s_{z1}+s_{z2})+J\vec s_1\cdot\vec s_2\right]t} = e^{-ib\,s_{z1}t}\;e^{-ib\,s_{z2}t}\;e^{-iJt\,\vec s_1\cdot\vec s_2}.
\end{equation}

Per il pezzo di campo, usando $s_{zi}=\sigma_{zi}/2$:

\begin{equation}
e^{-ib\,s_{zi}t} = e^{-i\frac{bt}{2}Z_i}.
\end{equation}

Confrontando con la convenzione Qiskit per la rotazione a un qubit, $R_z(\phi)=e^{-i\frac{\phi}{2}Z}$ (stessa struttura, stesso fattore $1/2$, già verificata per $R_{XX}$ nella Sezione~\ref{sec:UJ}), si identifica $\phi=bt$:

\begin{equation}
e^{-ib\,s_{zi}t} = R_z^{(i)}(bt).
\end{equation}

Il risultato finale, valido in qualunque ordine tra i tre fattori (poiché tutti commutano a due a due), è:

\begin{equation}
\boxed{\ket{\psi(t)} = R_z^{(1)}(bt)\,R_z^{(2)}(bt)\,U_J(t)\,\ket{\psi(0)}}
\end{equation}

con $U_J(t)$ come derivato nella Sezione~\ref{sec:UJ}. Questa espressione è \textbf{esatta}: per $D=0$ non è presente alcun errore di approssimazione, indipendentemente dal tempo $t$ simulato.

\section{Decomposizione esatta di $U_2(t) = e^{-iH_2t}$ (termine DM)}
\label{sec:U2}

\subsection{Riscrittura in operatori di Pauli}

Sostituendo $s_i^\alpha=\sigma_i^\alpha/2$ nel termine DM:

\begin{equation}
H_2 = D(s_{x1}s_{z2}-s_{z1}s_{x2}) = D\left(\frac{X_1}{2}\frac{Z_2}{2} - \frac{Z_1}{2}\frac{X_2}{2}\right) = \frac{D}{4}\left(X_1Z_2 - Z_1X_2\right).
\end{equation}

Si definiscono, per brevità di notazione, i due operatori a due qubit

\begin{equation}
P_1 \equiv X_1Z_2 \quad(\text{$X$ sul qubit 1, $Z$ sul qubit 2}), \qquad P_2 \equiv Z_1X_2 \quad(\text{$Z$ sul qubit 1, $X$ sul qubit 2}),
\end{equation}

cosicché $H_2 = \dfrac{D}{4}(P_1-P_2)$.

\subsection{Commutatività di $P_1$ e $P_2$}

Si calcolano esplicitamente i due prodotti, usando ancora $(A\otimes B)(C\otimes D)=(AC)\otimes(BD)$ e le relazioni $XZ=-iY$, $ZX=iY$:

\begin{align}
P_1P_2 &= (X_1Z_2)(Z_1X_2) = (X_1Z_1)\otimes(Z_2X_2) = (-iY_1)\otimes(iY_2) = (-i)(i)\,Y_1Y_2 = Y_1Y_2,\\
P_2P_1 &= (Z_1X_2)(X_1Z_2) = (Z_1X_1)\otimes(X_2Z_2) = (iY_1)\otimes(-iY_2) = (i)(-i)\,Y_1Y_2 = Y_1Y_2.
\end{align}

I due prodotti coincidono: $P_1P_2 = P_2P_1 = Y_1Y_2$, dunque $[P_1,P_2]=0$.

Ponendo $\theta \equiv \dfrac{Dt}{4}$, per l'argomento generale della Sezione~\ref{sec:BCH} si ha, esattamente e senza approssimazione:

\begin{equation}
e^{-i\theta(P_1-P_2)} = e^{-i\theta P_1}\,e^{+i\theta P_2}.
\label{eq:U2-factorized}
\end{equation}

\subsection{Circuito per i singoli termini: identità di Hadamard}

Si sfrutta l'identità $HZH=X$ (l'Hadamard scambia le basi computazionale e $X$), che si dimostra per calcolo diretto delle matrici $2\times2$:

\begin{equation}
H = \frac{1}{\sqrt2}\begin{pmatrix}1&1\\1&-1\end{pmatrix}, \qquad Z=\begin{pmatrix}1&0\\0&-1\end{pmatrix}.
\end{equation}

Calcolando prima $ZH$:

\begin{equation}
ZH = \begin{pmatrix}1&0\\0&-1\end{pmatrix}\frac{1}{\sqrt2}\begin{pmatrix}1&1\\1&-1\end{pmatrix} = \frac{1}{\sqrt2}\begin{pmatrix}1&1\\-1&1\end{pmatrix},
\end{equation}

e poi $H(ZH)$:

\begin{equation}
H(ZH) = \frac{1}{\sqrt2}\begin{pmatrix}1&1\\1&-1\end{pmatrix}\cdot\frac{1}{\sqrt2}\begin{pmatrix}1&1\\-1&1\end{pmatrix} = \frac{1}{2}\begin{pmatrix}1-1 & 1+1\\ 1+1 & 1-1\end{pmatrix} = \begin{pmatrix}0&1\\1&0\end{pmatrix} = X.
\end{equation}

Si conferma quindi $HZH=X$ (e, per la stessa identità applicata al contrario, $HXH=Z$, essendo $H^2=I$ e quindi $H$ auto-inversa).

Applicando l'Hadamard \emph{solo al qubit che porta $X$} nel termine considerato (qubit 1 per $P_1=X_1Z_2$):

\begin{equation}
X_1Z_2 = (H_1\otimes I_2)\,Z_1Z_2\,(H_1\otimes I_2),
\end{equation}

da cui, per l'esponenziale (l'Hadamard, essendo unitaria e auto-inversa, si porta dentro e fuori dall'esponenziale in modo standard, $e^{U A U^\dagger}=U e^A U^\dagger$):

\begin{equation}
e^{-i\theta\,X_1Z_2} = (H_1\otimes I_2)\;e^{-i\theta\,Z_1Z_2}\;(H_1\otimes I_2).
\end{equation}

Usando la convenzione Qiskit $R_{ZZ}(\phi)=e^{-i\frac{\phi}{2}ZZ}$ (Sezione~\ref{sec:UJ}), l'angolo da passare al gate è $\phi=2\theta$:

\begin{equation}
e^{-i\theta\,X_1Z_2} = (H\otimes I)\;R_{ZZ}(2\theta)\;(H\otimes I).
\label{eq:P1-circuit}
\end{equation}

Analogamente, per il secondo termine $P_2=Z_1X_2$, l'Hadamard va applicata al qubit 2 (che porta $X$):

\begin{equation}
Z_1X_2 = (I_1\otimes H_2)\,Z_1Z_2\,(I_1\otimes H_2) \quad\Longrightarrow\quad e^{+i\theta\,Z_1X_2} = (I\otimes H)\;R_{ZZ}(-2\theta)\;(I\otimes H),
\label{eq:P2-circuit}
\end{equation}

dove il segno dell'angolo si inverte rispetto all'Eq.~\eqref{eq:P1-circuit} a causa del segno $+i\theta$ nell'Eq.~\eqref{eq:U2-factorized}.

\subsection{Circuito completo per $U_2(t)$}

Combinando le Eq.~\eqref{eq:P1-circuit} e~\eqref{eq:P2-circuit} secondo l'ordine dell'Eq.~\eqref{eq:U2-factorized}, e sostituendo $\theta=Dt/4$ (quindi $2\theta=Dt/2$):

\begin{equation}
\boxed{U_2(t) = \Big[(I\otimes H)\,R_{ZZ}\!\Big(-\dfrac{Dt}{2}\Big)\,(I\otimes H)\Big]\cdot\Big[(H\otimes I)\,R_{ZZ}\!\Big(\dfrac{Dt}{2}\Big)\,(H\otimes I)\Big]}
\end{equation}

Anche questo blocco, isolato, è \textbf{esatto}: l'unica fonte di errore di approssimazione nell'intera simulazione è, come mostrato nella Sezione~\ref{sec:trotter-error}, l'accoppiamento tra $H_1$ e $H_2$.

\subsection{Nota sulla convenzione di segno del termine DM}

Se il termine DM fosse definito con segno opposto rispetto alla convenzione adottata, cioè $H_2' = D(s_{z1}s_{x2}-s_{x1}s_{z2}) = -H_2$, allora

\begin{equation}
U_2'(t) = e^{-iH_2't} = e^{+iH_2t} = U_2(-t).
\end{equation}

Concretamente, questo significa che \emph{entrambi} gli angoli dei due gate $R_{ZZ}$ cambiano segno simultaneamente: il blocco che nella convenzione adottata porta $R_{ZZ}(-Dt/2)$ diventerebbe $R_{ZZ}(+Dt/2)$ e viceversa. La struttura del circuito (sequenza $H$--$R_{ZZ}$--$H$, l'assegnazione di quale Hadamard va su quale qubit) resta invece identica: cambia unicamente il segno dei parametri $\theta$ passati ai gate $R_{ZZ}$, non la topologia del circuito.

\section{Perché $[A,B]=0$ implica fattorizzazione esatta}
\label{sec:BCH}

Tutte le fattorizzazioni derivate nelle Sezioni~\ref{sec:UJ}--\ref{sec:U2} si basano sullo stesso principio generale, che si dimostra qui una volta per tutte tramite la formula di Baker--Campbell--Hausdorff (BCH).

Per due operatori generici $A,B$ (non necessariamente commutanti), la formula BCH afferma:

\begin{equation}
e^{A}e^{B} = \exp\left(A+B+\frac{1}{2}[A,B]+\frac{1}{12}\big[A,[A,B]\big]+\frac{1}{12}\big[B,[B,A]\big]+\cdots\right),
\end{equation}

dove tutti i termini indicati con i puntini sono commutatori annidati di ordine sempre più alto, coinvolgenti esclusivamente $A$, $B$ e i loro commutatori successivi.

Se $[A,B]=0$, allora:
\begin{itemize}
\item il termine $\frac12[A,B]$ si annulla identicamente;
\item ogni commutatore annidato successivo, essendo costruito a partire da $[A,B]$ (o da commutatori che a loro volta si annullano), si annulla anch'esso identicamente, per induzione sulla profondità di annidamento.
\end{itemize}

Di conseguenza, \emph{l'intera serie infinita di correzioni} collassa a zero, non solo il primo termine:

\begin{equation}
[A,B]=0 \quad\Longrightarrow\quad e^{A}e^{B} = e^{A+B} \quad \text{esattamente}.
\end{equation}

Questo è il motivo per cui:
\begin{itemize}
\item la fattorizzazione di $U_J(t)$ nei tre fattori $R_{XX}, R_{YY}, R_{ZZ}$ (Sezione~\ref{sec:UJ}) è esatta, poiché $[XX,YY]=[YY,ZZ]=[ZZ,XX]=0$;
\item la fattorizzazione di $e^{-iH_1t}$ nel prodotto campo$\times$scambio (Sezione~\ref{sec:H1}) è esatta, poiché $[S^2,S_z]=0$;
\item la fattorizzazione di $U_2(t)$ nei due blocchi $P_1,P_2$ (Sezione~\ref{sec:U2}) è esatta, poiché $[P_1,P_2]=0$.
\end{itemize}

Il caso di $H_1$ e $H_2$ è invece concettualmente diverso: per $D\neq0$ si ha $[H_1,H_2]\neq0$ (il termine DM rompe la conservazione di $M$, come già verificato numericamente nel lavoro sul dimero VQE), la serie BCH \emph{non} si tronca, e approssimare $e^{-i(H_1+H_2)t}\approx e^{-iH_1t}e^{-iH_2t}$ significa scartare il termine $-\frac12[H_1,H_2]t^2$ e tutti i termini di ordine superiore della serie. Da qui la necessità del prodotto di Suzuki--Trotter, il cui errore viene quantificato nella sezione seguente.

\section{Derivazione dell'errore di Trotter}
\label{sec:trotter-error}

\subsection{Espansione in serie per un singolo passo}

Si consideri un singolo passo di Trotter di durata $\tau = t/N$, dove $N$ è il numero totale di passi per simulare il tempo totale $t$. Si confrontano l'evoluzione esatta e quella approssimata su questo singolo passo:

\begin{equation}
U_{\text{esatto}}(\tau) = e^{-i(H_1+H_2)\tau}, \qquad U_{\text{Trotter}}(\tau) = e^{-iH_1\tau}\,e^{-iH_2\tau}.
\end{equation}

Espandendo entrambe in serie di Taylor fino al second'ordine in $\tau$:

\begin{equation}
U_{\text{esatto}}(\tau) = I - i\tau(H_1+H_2) - \frac{\tau^2}{2}(H_1+H_2)^2 + O(\tau^3).
\end{equation}

Per il prodotto Trotter, si espande ciascun fattore separatamente,

\begin{equation}
e^{-iH_1\tau} = I - i\tau H_1 - \frac{\tau^2}{2}H_1^2 + O(\tau^3), \qquad e^{-iH_2\tau} = I - i\tau H_2 - \frac{\tau^2}{2}H_2^2 + O(\tau^3),
\end{equation}

e si moltiplicano, mantenendo tutti i termini fino a $O(\tau^2)$:

\begin{align}
U_{\text{Trotter}}(\tau) &= \left[I-i\tau H_1-\frac{\tau^2}{2}H_1^2\right]\left[I-i\tau H_2-\frac{\tau^2}{2}H_2^2\right]+O(\tau^3)\\
&= I - i\tau H_2 - \frac{\tau^2}{2}H_2^2 - i\tau H_1 + (-i\tau H_1)(-i\tau H_2) - \frac{\tau^2}{2}H_1^2 + O(\tau^3)\\
&= I - i\tau(H_1+H_2) - \frac{\tau^2}{2}\left(H_1^2+H_2^2\right) - \tau^2 H_1H_2 + O(\tau^3)\\
&= I - i\tau(H_1+H_2) - \frac{\tau^2}{2}\left(H_1^2+H_2^2+2H_1H_2\right) + O(\tau^3).
\end{align}

\subsection{Differenza tra le due espansioni}

Si espande ora $(H_1+H_2)^2$ nell'espressione esatta:

\begin{equation}
(H_1+H_2)^2 = H_1^2 + H_1H_2 + H_2H_1 + H_2^2.
\end{equation}

Sottraendo le due espansioni termine a termine:

\begin{align}
U_{\text{esatto}}(\tau) - U_{\text{Trotter}}(\tau) &= -\frac{\tau^2}{2}\left(H_1^2+H_1H_2+H_2H_1+H_2^2\right) + \frac{\tau^2}{2}\left(H_1^2+H_2^2+2H_1H_2\right) + O(\tau^3)\\
&= -\frac{\tau^2}{2}\left(H_1H_2+H_2H_1\right) + \frac{\tau^2}{2}\left(2H_1H_2\right) + O(\tau^3)\\
&= \frac{\tau^2}{2}\left(2H_1H_2-H_1H_2-H_2H_1\right) + O(\tau^3)\\
&= \frac{\tau^2}{2}\left(H_1H_2-H_2H_1\right) + O(\tau^3).
\end{align}

Riconoscendo $H_1H_2-H_2H_1=[H_1,H_2]$, si ottiene il risultato:

\begin{equation}
\boxed{U_{\text{esatto}}(\tau) - U_{\text{Trotter}}(\tau) = \frac{\tau^2}{2}\,[H_1,H_2] + O(\tau^3)}
\end{equation}

L'errore per singolo passo è dunque $O(\tau^2)=O\big((t/N)^2\big)$, in accordo con quanto riportato negli appunti del relatore, ed è proporzionale al commutatore $[H_1,H_2]$: coerentemente, per $D=0$ (dove $[H_1,H_2]=0$, si veda la Sezione~\ref{sec:H1}) l'errore si annulla identicamente, indipendentemente da $\tau$.

\subsection{Errore totale su $N$ passi}

L'evoluzione temporale completa fino al tempo $t$ è approssimata dal prodotto di $N$ passi identici:

\begin{equation}
U_{\text{Trotter}}(\tau)^N \approx U_{\text{esatto}}(t), \qquad \tau = t/N.
\end{equation}

Per stimare l'errore totale, si usa la disuguaglianza triangolare applicata alla norma operatoriale $\|\cdot\|$ (ad esempio la norma spettrale), che permette di limitare l'errore accumulato su $N$ passi con l'errore per singolo passo moltiplicato per $N$:

\begin{equation}
\left\|U_{\text{esatto}}(t) - U_{\text{Trotter}}(\tau)^N\right\| \;\lesssim\; N \cdot \left\|U_{\text{esatto}}(\tau)-U_{\text{Trotter}}(\tau)\right\| = N\cdot O(\tau^2).
\end{equation}

Sostituendo $\tau=t/N$:

\begin{equation}
\left\|U_{\text{esatto}}(t) - U_{\text{Trotter}}(\tau)^N\right\| \;\lesssim\; N\cdot O\!\left(\frac{t^2}{N^2}\right) = O\!\left(\frac{t^2}{N}\right).
\end{equation}

L'errore \textbf{totale} decresce quindi come $1/N$: si tratta di un metodo del prim'ordine. Raddoppiando il numero di passi $N$, a parità di tempo totale simulato $t$, l'errore totale si dimezza.

\subsection{Scaling rispetto al tempo totale simulato}

Dallo scaling $O(t^2/N)$ dell'errore totale derivato sopra, si ricavano due conseguenze pratiche:

\begin{itemize}
\item \textbf{A $N$ fisso}: raddoppiando il tempo totale di simulazione $t\to 2t$, l'errore scala come $(2t)^2/N = 4\cdot(t^2/N)$, ossia \textbf{quadruplica}.
\item \textbf{A errore fisso}: per mantenere lo stesso errore totale mentre si raddoppia $t$, occorre che $t^2/N$ resti costante. Se $t\to2t$, allora $(2t)^2/N' = t^2/N$ implica $N' = 4N$: il numero di passi deve quadruplicare.
\end{itemize}

In generale, per mantenere l'errore costante al variare di $t$, il numero di passi deve scalare come $N\propto t^2$.

\paragraph{Implicazione pratica.} Poiché il numero di gate nel circuito di Trotter è proporzionale al numero di passi $N$, il costo circuitale (profondità, numero di gate a due qubit) cresce \textbf{quadraticamente} con il tempo totale simulato, a errore fissato. Questo è uno dei limiti pratici noti del prodotto di Trotter al prim'ordine, e motiva l'uso, in letteratura, di schemi di ordine superiore (ad esempio il prodotto di Trotter simmetrico, o formule di Suzuki di ordine 2, con errore per passo $O(\tau^3)$ anziché $O(\tau^2)$) quando è necessario simulare tempi lunghi. Per il presente task, coerentemente con quanto richiesto dal relatore, ci si limita al prim'ordine.

\subsection{Metodi di verifica numerica}

La derivazione teorica dello scaling dell'errore può essere verificata numericamente con due approcci standard, entrambi compatibili con gli strumenti già impiegati nel resto del progetto:

\begin{enumerate}
\item \textbf{Norma operatoriale.} Si costruisce esplicitamente l'evoluzione esatta $U_{\text{esatto}}(t)=\exp(-iHt)$ (ad esempio tramite \texttt{scipy.linalg.expm}, oppure tramite gli strumenti di algebra lineare di Qiskit come \texttt{Operator}) e l'evoluzione di Trotter $U_{\text{Trotter}}(\tau)^N$ come prodotto esplicito di matrici. Si calcola quindi $\left\|U_{\text{esatto}}-U_{\text{Trotter}}^N\right\|$ (norma spettrale o di Frobenius) al variare di $N$, verificando che lo scaling numerico riproduca l'andamento teorico $\propto 1/N$.
\item \textbf{Fidelity sullo stato evoluto.} Si applicano entrambe le evoluzioni a un medesimo stato iniziale $\ket{\psi_0}$, ottenendo $\ket{\psi_{\text{esatto}}(t)}$ e $\ket{\psi_{\text{Trotter}}(t)}$, e si calcola la fidelity
\begin{equation}
\mathcal F(t) = \left|\langle \psi_{\text{esatto}}(t) \,|\, \psi_{\text{Trotter}}(t) \rangle\right|,
\end{equation}
verificando che $1-\mathcal F(t) \propto t^2/N$ (a $t$ fissato, al variare di $N$). Questa è la stessa metodologia di fidelity già impiegata nella validazione del VQE nel resto del progetto, e si presta naturalmente al riutilizzo del codice esistente.
\end{enumerate}

\section{Riepilogo dei risultati}

La Tabella~\ref{tab:riepilogo} riassume tutti i blocchi derivati, distinguendo quali sono esatti e quale introduce l'unico errore di approssimazione dell'intera costruzione.

\begin{table}[h]
\centering
\begin{tabular}{@{}p{3.3cm}p{7.7cm}c@{}}
\toprule
\textbf{Blocco} & \textbf{Espressione} & \textbf{Natura} \\
\midrule
$U_J(t)$ (scambio) & $R_{XX}(Jt/2)\,R_{YY}(Jt/2)\,R_{ZZ}(Jt/2)$ & Esatta \\[4pt]
$e^{-iH_1t}$ ($D=0$) & $R_z^{(1)}(bt)\,R_z^{(2)}(bt)\,U_J(t)$ & Esatta \\[4pt]
$U_2(t)$ (termine DM) & $\big[(I\otimes H)R_{ZZ}(-Dt/2)(I\otimes H)\big]\cdot\big[(H\otimes I)R_{ZZ}(Dt/2)(H\otimes I)\big]$ & Esatta \\[4pt]
$e^{-i(H_1+H_2)t}$, $D\neq0$ & $\left(e^{-iH_1t/N}\,e^{-iH_2t/N}\right)^N$ & Approssimata, errore $O(t^2/N)$ \\
\bottomrule
\end{tabular}
\caption{Riepilogo dei blocchi di evoluzione temporale derivati per l'Hamiltoniana del dimero.}
\label{tab:riepilogo}
\end{table}

\end{document}
